\documentclass[12pt]{article}
\usepackage{amsmath, amssymb}
\usepackage{graphicx}
\usepackage{float}
\usepackage[letterpaper, margin=1in]{geometry}
\usepackage[colorlinks=true, linkcolor=blue, urlcolor=blue]{hyperref}
\usepackage{listings}
\usepackage{booktabs}
\usepackage{siunitx}
\usepackage{booktabs, xcolor, colortbl, array, caption}
\usepackage{lettrine}
\usepackage{tikz}
\begin{document}
\pagestyle{empty}
\begin{center}
    \Huge \textbf{Conclusion}
\end{center}

Both the t-test and t-interval imply that there is no real difference in Shohei Ohtani's batting ability when pitching vs not pitching. Thus, Ohtani should pitch whenever he can. There is likely no reason to prevent him from pitching in order to improve his batting performance or vice versa. However, an experiment would be needed to be run to confirm this result, as cause and effect cannot be inferred, because this is an observational study rather than an experiment.

This is one limitation of my conclusions. Another is that the sample size for pitching was much smaller than when not pitching, which increases the variability and the chance of a type II error. I did not control for weather, opponent, rest days, injuries, etc., which could effect the conclusion. Game-to-game wOBA is noisy, and I did not control for the number of plate appearances, besides only using games with at least 3. A more accurate way to run the tests would be to weight the wOBA by plate appearances, although the math would get much more complicated and would likely produce similar results.

The study could be improved by increasing the sample size, meaning it would have to be repeated in the future after Ohtani had played more games. It could also be improved by accounting for variables like weather and rest, but a more advanced statistical method, such as regression, would need to be used, rather than a t-test.

Overall, the results suggest that Shohei's batting is consistent regardless of whether or not he is pitching. While there is no statistically significant difference, the data and methods I used have limitations, meaning a more complete and extensive study may be able to identify a difference. Nevertheless, the findings highlight how special of a player Shohei is, as he is able to hit at an elite level, even while simultaneously performing of the most physically demanding roles in all of sports.
\end{document}